\documentclass[11pt,a4paper]{article}

% ---------------------------------------------------------
% PREAMBLE (stable, warning-free, Overleaf-safe)
% ---------------------------------------------------------
\usepackage[utf8]{inputenc}
\usepackage[T1]{fontenc}
\usepackage{lmodern}
\usepackage{amsmath, amssymb, amsfonts}
\usepackage{bm}
\usepackage{geometry}
\usepackage{graphicx}
\usepackage{hyperref}
\usepackage{cite}
\usepackage{authblk}
\usepackage{physics}
\usepackage{mathtools}

\geometry{margin=2.4cm}

\title{\textbf{Companion LXVI: Ensemble Variance of Persistence Diagrams in the Relaxation-Driven Cyclic Cosmology}}
\author{Michael Lehmann \\ \texttt{mi.lehmann@gmx.de}}
\date{April 2026}

\begin{document}
\maketitle

\begin{abstract}
The ensemble variance of persistence diagrams quantifies the fluctuations of 
topological structure across realizations of a cosmological field. In weak-lensing 
analyses, this variance measures the stability of connected components, loops, and 
void-like structures under cosmic variance, noise, and nonlinear evolution. In the 
standard cosmological model, the variance reflects the nonlinear matter 
distribution and the projection of large-scale structure. In the Relaxation-Driven 
Cyclic Cosmology (RDCC), the scale-dependent evolution of the gravitational 
potential modifies the fluctuation spectrum of topological features in a 
characteristic way. This Companion develops a continuous description of ensemble 
variance in RDCC, deriving how the relaxation scale influences topological 
fluctuations, persistence-lifetime scatter, and the observational signatures in 
weak-lensing surveys. The resulting framework provides a distinctive probe of the 
RDCC gravitational field through topological variability.
\end{abstract}
% ---------------------------------------------------------
\section{Introduction}

Persistence diagrams encode the birth and death of topological features in 
weak-lensing convergence fields. While Fréchet means summarize the average 
topology of an ensemble, the ensemble variance quantifies the fluctuations around 
this mean. This variance is sensitive to nonlinear structure formation, cosmic 
variance, and the stability of topological features.

In the standard cosmological model, ensemble variance reflects the nonlinear matter 
distribution and the projection of large-scale structure. In RDCC, the 
gravitational potential evolves differently. The relaxation mechanism suppresses 
the decay of small-scale modes and enhances the coherence of large-scale modes. 
This modifies the fluctuation spectrum of topological features in a scale-dependent 
manner, producing a distinctive signature in ensemble-variance analyses.
% ---------------------------------------------------------
\section{Ensemble Variance of Persistence Diagrams}

Given an ensemble of persistence diagrams $\{D_{i}\}$ with Fréchet mean $D^{*}$, 
the ensemble variance is defined as

\begin{equation}
    \mathrm{Var}(D)
    = \frac{1}{N}
      \sum_{i=1}^{N}
      W_{p}(D_{i}, D^{*})^{2},
\end{equation}

where $W_{p}$ is the $p$-Wasserstein distance.

This variance measures the spread of topological features across the ensemble.
% ---------------------------------------------------------
\section{RDCC Modification of Persistence Diagrams}

In RDCC, the convergence evolves as
\begin{equation}
    \kappa_{\mathrm{RDCC}}(\ell)
    = \kappa(\ell)
      \left[
        1 - \chi_{\kappa}
        \left(\frac{\ell}{\ell_{\mathrm{rel}}}\right)^{\alpha}
      \right].
\end{equation}

This modifies the birth and death thresholds of topological features:

\begin{align}
    b_{\mathrm{RDCC}} &= b \left[ 1 - \chi_{b} (\ell/\ell_{\mathrm{rel}})^{\alpha} \right], \\
    d_{\mathrm{RDCC}} &= d \left[ 1 - \chi_{d} (\ell/\ell_{\mathrm{rel}})^{\alpha} \right].
\end{align}

The ensemble of RDCC diagrams becomes
\begin{equation}
    \{ D_{i,\mathrm{RDCC}} \}
    = \left\{
        (b_{i,\mathrm{RDCC}}, d_{i,\mathrm{RDCC}})
      \right\}.
\end{equation}
% ---------------------------------------------------------
\section{Ensemble Variance in RDCC}

The RDCC ensemble variance is

\begin{equation}
    \mathrm{Var}_{\mathrm{RDCC}}(D)
    = \frac{1}{N}
      \sum_{i=1}^{N}
      W_{p}(D_{i,\mathrm{RDCC}}, D^{*}_{\mathrm{RDCC}})^{2}.
\end{equation}

To first order in the relaxation parameter,

\begin{equation}
    \mathrm{Var}_{\mathrm{RDCC}}(D)
    = \mathrm{Var}(D)
      \left[
        1 - \chi_{\mathrm{Var}}
        \left(\frac{\ell}{\ell_{\mathrm{rel}}}\right)^{\alpha}
      \right].
\end{equation}

This modification reflects the relaxation-driven change in topological 
fluctuations.
% ---------------------------------------------------------
\section{Topological Fluctuation Signatures of RDCC}

The relaxation mechanism enhances the coherence of large-scale modes. Ensemble 
variance therefore exhibits:

\begin{itemize}
    \item reduced variance at large scales (enhanced coherence),
    \item enhanced variance at small scales (suppressed small-scale potential),
    \item a characteristic turnover near $\ell_{\mathrm{rel}}$,
    \item modified scaling with smoothing scale and redshift,
    \item altered clustering of diagrams around the Fréchet mean.
\end{itemize}

These signatures provide a direct probe of the RDCC gravitational field through 
topological fluctuations.
% ---------------------------------------------------------
\section{Conclusion}

The ensemble variance of persistence diagrams in the Relaxation-Driven Cyclic 
Cosmology reflects the scale-dependent evolution of the gravitational potential and 
the nonlinear matter distribution. The relaxation mechanism modifies the 
fluctuation spectrum of topological features in a scale-dependent manner, 
producing distinctive signatures in weak-lensing surveys. These effects provide a 
direct observational window into the relaxation scale and the temporal structure of 
the RDCC gravitational field. The present Companion establishes the theoretical 
foundation for future studies of topological variability, nonlinear structure, and 
the interplay between light and gravity in RDCC.

\bibliographystyle{apsrev4-2}
\nocite{*}
\bibliography{references}


\end{document}
